% ---------------------------------------------------------
\section{Continuity}

% ---------------------------------------------------------
\section*{Definition}

\begin{propositionbox}{Definition (Continuous Function)}
Let $f$ be a function whose domain contains an open interval $I$.
\begin{enumerate}
  \item $f$ is \emph{continuous at a point} $c \in I$ if
        $\displaystyle\lim_{x \to c} f(x) = f(c)$.
  \item $f$ is \emph{continuous on} $I$ if $f$ is continuous at
        every $c \in I$. If $f$ is continuous on $(-\infty, \infty)$,
        $f$ is said to be \emph{continuous everywhere}.
\end{enumerate}
\end{propositionbox}

\begin{remark*}[Three-condition test]
$f$ is continuous at $c$ if and only if all three of the following hold:
\begin{enumerate}
  \item $\displaystyle\lim_{x \to c} f(x)$ exists,
  \item $f(c)$ is defined, and
  \item $\displaystyle\lim_{x \to c} f(x) = f(c)$.
\end{enumerate}
The three conditions make explicit what can go wrong: the limit may
fail to exist, the function may be undefined at $c$, or the two
values may disagree.
\end{remark*}

% ---------------------------------------------------------
\section*{Continuity on Closed Intervals}

\begin{propositionbox}{Definition (Continuity on Closed Intervals)}
Let $f$ be defined on the closed interval $[a, b]$.
$f$ is \emph{continuous on} $[a, b]$ if:
\begin{enumerate}
  \item $f$ is continuous on $(a, b)$,
  \item $\displaystyle\lim_{x \to a^+} f(x) = f(a)$, and
  \item $\displaystyle\lim_{x \to b^-} f(x) = f(b)$.
\end{enumerate}
\end{propositionbox}

\begin{remark*}
Condition (2) is called \emph{continuity from the right} at $a$;
condition (3) is \emph{continuity from the left} at $b$.
The same idea applies to half-open intervals $[a, b)$ and $(a, b]$.
\end{remark*}

% ---------------------------------------------------------
\section*{Classification of Discontinuities}

\begin{propositionbox}{Definition (Types of Discontinuity)}
When $f$ fails to be continuous at $a$, the discontinuity falls into
one of three categories.
\begin{enumerate}
  \item \textbf{Removable discontinuity.}
        $\lim_{x \to a} f(x) = L$ exists, but either $L \neq f(a)$
        or $f(a)$ is undefined.
        The discontinuity is ``removable'' because redefining
        $f(a) = L$ produces a function continuous at $a$ that agrees
        with $f$ everywhere else.

  \item \textbf{Jump discontinuity.}
        $\lim_{x \to a^+} f(x) = L$ and $\lim_{x \to a^-} f(x) = M$
        both exist, but $L \neq M$.
        The two-sided limit does not exist; the graph visually
        ``jumps'' from one value to another as $x$ crosses $a$.

  \item \textbf{Infinite discontinuity.}
        $f$ is unbounded near $x = a$: the value of $f$ becomes
        arbitrarily large (or large and negative) as $x$ approaches
        $a$.
        At least one one-sided limit is $\pm\infty$, and $x = a$
        is a vertical asymptote of $f$.
\end{enumerate}
\end{propositionbox}

\begin{remark*}[Canonical examples]
Removable: $f(x) = \dfrac{x^2 - 1}{x - 1}$ at $x = 1$ --- the
limit is $2$ but $f(1)$ is undefined.
Jump: $f(x) = \lfloor x \rfloor$ at every integer --- the left- and
right-hand limits differ by $1$.
Infinite: $f(x) = \dfrac{1}{x}$ at $x = 0$ --- the function grows
without bound from both sides.
\end{remark*}

% ---------------------------------------------------------
\section*{Properties of Continuous Functions}

\begin{theorembox}{Theorem (Properties of Continuous Functions)}
Let $f$ and $g$ be continuous on an interval $I$, let $c \in \mathbb{R}$,
and let $n$ be a positive integer.
The following functions are continuous on $I$:

\medskip
\begin{tabular}{@{}p{0.22\textwidth}p{0.70\textwidth}@{}}
\textbf{Sums/Differences}  & $f \pm g$ \\[6pt]
\textbf{Constant Multiple} & $c \cdot f$ \\[6pt]
\textbf{Products}          & $f \cdot g$ \\[6pt]
\textbf{Quotients}         & $f/g$ \quad (provided $g \neq 0$ on $I$) \\[6pt]
\textbf{Powers}            & $f^n$ \\[6pt]
\textbf{Roots}             & $\sqrt[n]{f}$ \quad (if $n$ is even, require $f \geq 0$ on $I$) \\[6pt]
\textbf{Compositions}      & If $f$ is continuous on $I$ with range $J$, and $g$ is
                             continuous on $J$, then $g \circ f$ is continuous on $I$.
\end{tabular}
\end{theorembox}

\begin{remark*}
Polynomials are continuous everywhere.
Rational functions are continuous on any interval that avoids the
zeros of the denominator.
\end{remark*}

% ---------------------------------------------------------
\section*{Continuous Functions}

\begin{theorembox}{Theorem (Continuous Functions)}
The following functions are continuous on their domains:

\medskip
\begin{tabular}{@{}p{0.45\textwidth}p{0.45\textwidth}@{}}
$\sin(x)$     & $\cos(x)$           \\[4pt]
$\tan(x)$     & $\cot(x)$           \\[4pt]
$\sec(x)$     & $\csc(x)$           \\[4pt]
$\ln(x)$      & $a^x \quad (a > 0)$ \\[4pt]
$\sqrt[n]{x}$ & ($n$ a positive integer) \\
\end{tabular}
\end{theorembox}

% ---------------------------------------------------------
\section*{The Intermediate Value Theorem}

\begin{theorembox}{Theorem (Intermediate Value Theorem)}
Let $f$ be continuous on $[a, b]$ and, without loss of generality,
let $f(a) < f(b)$.
Then for every value $y$ with $f(a) < y < f(b)$, there exists at
least one $c \in (a, b)$ such that $f(c) = y$.
\end{theorembox}

\begin{remark*}[Intuition]
A continuous function cannot skip values.
If $f$ is at height $f(a)$ on the left and height $f(b)$ on the right,
it must pass through every height in between at least once.
\end{remark*}

\begin{remark*}[Root-finding application]
The IVT is the theoretical foundation of the \emph{Bisection Method}.
If $f(a) < 0$ and $f(b) > 0$ with $f$ continuous on $[a, b]$, then
there is at least one root $c \in (a, b)$ where $f(c) = 0$.
Evaluating $f$ at the midpoint $d = (a + b)/2$ and replacing whichever
endpoint has the same sign as $f(d)$ halves the interval at each step,
converging to the root.
\end{remark*}

\begin{remark*}[What the IVT does not say]
The theorem guarantees existence of at least one $c$, but gives no
information about how many such values exist or where they are.
It also requires continuity on the \emph{entire} closed interval
$[a, b]$; a single discontinuity is enough to invalidate the
conclusion.
\end{remark*}